% Section 9.6, from the summary table of lectures/L09/appendix/a_theory.md and from
% lectures/L09/README.md: the objectives, the evaluation questions and the next lesson.

\section{Sammanfattning}
\label{c9:sec:sammanfattning}

\begin{booktable}{@{}p{12.5em}L@{}}
  \thead{Begrepp} & \thead{Formel} \\
  \midrule
  Grader till radianer & $v_{\text{rad}} = v_{\text{grad}} \cdot \pi/180^{\circ}$ \\
  Vinkelhastighet & $\omega = 2\pi f = 2\pi/T$ \\
  Periodtid & $T = 1/f$ \\
  Växelspänningsekvation & $u(t) = \abs{U}\sin(\omega t + \delta)$ \\
\end{booktable}

\needspace{6\baselineskip}
\subsection*{Det här ska du kunna}
Efter det här kapitlet ska du:
\begin{itemize}
  \item Känna till de trigonometriska funktionerna och deras parametrar.
  \item Kunna använda vinkelmåtten grader och radianer och omvandla mellan dem.
\end{itemize}

\testadigsjalv
\begin{enumerate}
  \item Omvandla $150^{\circ}$ till radianer och $\dfrac{5\pi}{6}$ till grader.
  \item En växelspänning har amplituden $5\,\text{V}$, frekvensen $50\,\text{Hz}$ och fasen
    $\pi/4\,\text{rad}$. Skriv spänningens ekvation $u(t)$.
\end{enumerate}

\needspace{6\baselineskip}
\subsection*{Nästa kapitel}
\Kapref{10} handlar om exponentialfunktioner och logaritmer samt decibel.
